\renewcommand{\thetable}{3.A.11}
\begin{tablalafrox}{Consultas abiertas. Fuente: contraste de las Bases y de los Subdocumentos 1 y 2.}%
  {tab:3-A-11}%
  {F{0.060} Y F{0.164} F{0.219} F{0.167}}%
  {\cab{ID} & \cab{Pregunta} & \cab{Efecto en el alcance} & \cab{Supuesto de oferta} & \cab{Responsable y momento}}
  V-01 & ¿Cuántas instalaciones debe cubrir la solución: cinco según las entrevistas o seis según la volumetría del caso? & Número de sitios con borde operacional y de olas. & S-22: se parametriza por sitio. & Gerencia de Operaciones, mes 1. \\
  V-02 & La ponderación del Formulario T-21 suma 98 \%; ¿cuál es el peso correcto? & Ninguno sobre el alcance técnico. & Se usan los pesos publicados. & Consulta al mandante. \\
  V-03 & Decisión de oferta adoptada: cinco ambientes. Sigue pendiente la aclaración formal de la diferencia entre \textit{Bases Administrativas} (2026, art. 24, p. 17; Formulario E-25, hito H3, p. 74), y \textit{Bases Técnicas Transversales} (2026, cap. 4, RT-04.01, p. 10). & Cinco ambientes operativos para H3. & D-11: DEV, QA, PREPROD, PROD y DR, conforme a la infraestructura declarada en el Subdocumento 1 (LafroX, 2026b, sección «Presentación de la empresa») y \textit{Bases Técnicas Transversales} (2026, cap. 4, RT-04.01, p. 10). No acredita respuesta del mandante. & Consulta formal al mandante pendiente; la decisión de oferta está adoptada. \\
  V-04 & Resuelta documentalmente: los códigos RT se identifican por documento de origen; no son equivalentes por compartir número. & Correspondencia de obligaciones en el T-12, sin sustituir unas por otras. & La oferta exige sincronización de reparto en máximo 10 minutos y de CD en máximo 2 horas, conforme a D-08. \textit{Bases Técnicas Transversales} (2026, cap. 3, RT-03.12, p. 9): sincronización automática, reconciliación determinista y bitácora (RNG-12); \textit{Bases Técnicas Transversales} (2026, cap. 3, RT-03.13, p. 9): declarar funciones no disponibles sin conexión y procedimiento que las suple. & No requiere respuesta para establecer la correspondencia documental; su cumplimiento se verifica en la oferta. \\
  V-05 & Resuelta documentalmente: los cortes relatados no reducen la autonomía mínima exigida. & Operación desconectada de terreno y CD. & S-24: 14 h sin cobertura en terreno y 24 h sin enlace en CD (Caso 02, cap. 15, RT-03.10; LafroX, 2026a, Anexo 2.1, RNF-RESIL; D-08). & No requiere respuesta; cumplimiento pendiente de demostrar mediante pruebas. \\
  V-06 & ¿Qué mecanismos de registro, verificación de identidad y recuperación de acceso se usarán por perfil sin correo corporativo? & Separar personal interno y personas externas; \textit{Bases Técnicas Transversales} (2026, cap. 12, RT-12.12, p. 26), exige acceso autoservido y seguro para estas últimas. & El RUT y la verificación por supervisor son una propuesta de identificación, no un mecanismo completo de recuperación. El Caso 02 (cap. 15, RT-12.12) identifica como personas usuarias externas a clientes del canal tradicional y food service, cadenas de supermercados, empresas transportistas y sus conductores, y proveedores. & Jefe de TI, mes 2; mecanismos y canales por confirmar. \\
  V-07 & ¿Qué proveedores entregan sin lote visible y qué fuente documental permite conocerlo? ¿Qué tratamiento se aprueba cuando el lote no puede determinarse? & Captura estructurada en recepción y tratamiento de productos con lote desconocido. & Ingreso manual solo cuando el lote del proveedor es conocido por una fuente verificable; RF-01.02 impide continuar con el campo vacío para SKU con trazabilidad obligatoria. No se inventa ni sustituye el lote (LafroX, 2026a, Anexo 2.2, S-02; RF-01.06). & Jefa de Calidad, mes 2; política ante lote desconocido pendiente de aprobación. \\
  V-08 & ¿Existen antecedentes suficientes y confiables para determinar el saldo inicial de envases por cliente y transportista? & Carga inicial conciliada de M8. & Las anotaciones en cuadernos no demuestran ausencia de antecedentes. RNG-07 y (LafroX, 2026a, Anexo 2.2, S-10), requieren control por saldo. El conteo inicial es una alternativa cuya necesidad y momento deben confirmarse. & Gerencia de Operaciones, ola de reparto; saldos y conciliación por validar con Finanzas. \\
  V-09 & ¿Qué extensión, calidad y formato tiene la historia de ventas y pedidos disponible en el ERP para reposición y migración? & Datos de RF-10.01 y alcance de migración. & S-25. El Caso 02 (cap. 15, RT-05.15) exige migrar tres años de ventas y pedidos; esta exigencia no acredita su disponibilidad ni su calidad efectiva en el ERP. & Jefe de abastecimiento, mes 3; disponibilidad y calidad por confirmar con TI. \\
  V-10 & ¿Qué registros históricos permiten conocer la hora real de entrega y el monto esperado de rendición, y con qué confiabilidad? & Línea base histórica de OTIF y cuadratura, separada de la medición futura. & Desde marcha blanca, RF-06.06 captura la hora y RF-07.01 el monto esperado. Esto no reconstruye la historia anterior. OTIF conserva la definición de (LafroX, 2026a, Anexo 2.2, S-01), y RNG-09; antecedentes y período de línea base pendientes de validación. & Hora de entrega: Gerencias Comercial y de Operaciones (S-01); monto de rendición: Gerencia de Finanzas; mes 1. \\
  V-11 & ¿Qué umbrales y duraciones de excursión térmica aplica cada tipo de producto? & Parámetros pendientes de RNG-04. & El Anexo 2.2 del Subdocumento 2 (LafroX, 2026a, S-04) y RNG-04 ya definen alerta para excursión menor y transitoria, y bloqueo preventivo para crítica y sostenida; Calidad decide liberar, bloquear o rechazar. Los valores por producto siguen pendientes. Hasta su aprobación, toda lectura fuera del rango de almacenamiento declarado en la ficha del SKU se trata preventivamente como crítica; la marcha blanca no se cierra sin la matriz de parámetros aprobada por Calidad. & Jefa de Calidad, mes 2; matriz de parámetros por aprobar. \\
  V-12 & ¿Cuál es la fecha de inicio del contrato? & Calendario real de los meses 16 y 21. & S-17. & Contraparte Técnica, firma del contrato. \\
  V-13 & ¿Qué evidencia de trazabilidad acepta el proveedor de lácteos para restablecer la relación comercial tras la suspensión cuyo vencimiento estaba previsto en septiembre de 2026 según el caso? & Prioridad de RF-01.02 y RF-09.01 en la Etapa 1. & La trazabilidad de lote va primero en la secuencia de la Etapa 1. & Jefa de Calidad, mes 1. \\
  V-14 & ¿Cuál es la fecha efectiva de retiro del planificador y su disponibilidad para transferir y validar las reglas? & Calendario de talleres y validaciones de S-16. & (LafroX, 2026a, Anexo 2.2, S-16), compromete transferencia temprana; LafroX propone talleres en meses 1 a 3. Ante retiro anticipado, se valida con el ayudante y datos históricos; no confirma la fecha efectiva. & Gerencia de Operaciones y planificador, mes 1. \\
  V-15 & ¿En qué condiciones contractuales se accede a la telemetría del tercero? & Integración de M12. & S-19. & Jefe de TI, mes 2. \\
  V-16 & ¿Cuáles son las interfaces disponibles del ERP? & Integración y migración. & S-18. & Jefe de TI, mes 1. \\
  V-17 & Si la fecha de inicio ubica el mes 16 o el mes 21 en septiembre o diciembre, ¿cómo se compatibilizará formalmente el calendario contractual de \textit{Bases Administrativas} (2026, art. 17.1, p. 12) con las restricciones operacionales que prohíben pasos a producción en septiembre y diciembre? & Calendario de los pasos a producción. & S-17: inicio en un mes sin conflicto. & Mandante, mediante consulta escrita por el canal oficial durante el período de consultas de la licitación; cierre del asunto mediante respuesta formal. \\
  V-18 & ¿Qué mecanismo acepta el CLIENTE para el acuse tributario de receptores sin firma electrónica avanzada y cómo permite el ERP consultar y conciliar su estado? & Validación del mecanismo de F-09, RF-06.14 y RF-12.13; se conserva el ERP como emisor. & Integración comprometida: RF-06.14 transmite evidencia al ERP y concilia el estado del acuse con la guía. (LafroX, 2026a, Anexo 2.2, S-15), define alternativas de evidencia de entrega; no acredita aceptación del mecanismo de acuse ni capacidades verificadas del ERP. & Jefe de TI y Gerencia de Finanzas, mes 1; mecanismo y consulta de estado por confirmar. \\
\end{tablalafrox}
